% =====================================================================
% c01.tex —— 第 2 节末（p006 上部）+ 第 3 节「導来圏」（p006–p009）
% 底本：植田一石《ホモロジー的ミラー対称性》讲义录
% 公式一律以 work/pages/pNNN.png 页面图像为准逐符号核对。
%
% [Tutor Note: 本 chunk 起首承接 c00 末句（p005 页末「Kähler 多様体とその上の正則関数の組から」），自「決まる作用汎関数によって…」接续，非新段落。]
% [Tutor Note: BRIEF 记 c01 有 9 个编号公式，该数来自文字层命中（含正文对同一编号的交叉引用）。按页面图像核对，p006–p009 上实际编号公式为 (2.7)、(3.1)–(3.5) 共 6 条；此处按实际条数标注，未凑数、未给无编号公式加 tag。]
% =====================================================================

\noindent 所确定的作用泛函所指定、具有 $(2,2)$ 超对称性的 2 维场论\footnote{关于作用泛函的具体表达式，例如参见 [HKK$^{+}$03, 13 章]。}\footnote{Landau–Ginzburg 模型与以奇点论为中心的代数几何密切相关，这一点在 1980 年代末被认识到 [Mar90, VV89]。}。[GP90] 的推广 [Bat94, BD96, BB96] 中环面几何走到前台，Landau–Ginzburg 模型并不显式登场；但通过与外围的环面 Fano 流形的镜像对称的关系，Landau–Ginzburg 模型再次露出身影\footnote{例如参见 [Wit93] 与 [Giv95]。另外，猜想 2.5 中出现的洛朗多项式也与代数环面 $(\C^{\times})^n$ 组合，给出 Landau–Ginzburg 模型。}。另一方面，[GP90] 的另一个推广 [BH93] 中，Landau–Ginzburg 模型起中心作用。这在历史上与基于环面几何的镜像流形构造相比不太受关注，但受 Gromov–Witten 理论在 Landau–Ginzburg 模型中的类似物\footnote{[FJR13, PV] 及其参考文献。}的发展推动，近年被活跃地研究\footnote{基于环面几何的镜像流形构造与基于可逆多项式的镜像流形构造的关系是有趣的问题。试图统一理解它们的工作有 [Bor13] 与 [Cla] 等。}。此情形下的 Landau–Ginzburg 模型由可逆多项式（invertible polynomial）\footnote{可逆多项式与 Delsarte 多项式 [Del95, Shi86] 实质上是同一概念。}给出，同调镜像对称的主张大致如下：

\para{猜想 2.7.} 对任意可逆多项式 $f$，记其转置为 $\check{f}$，则存在如下增广三角范畴的等价：
\begin{equation}
    D^b\mathfrak{Fuk}\,f\cong\operatorname{hmf}(\check{f}). \tag{2.7}
\end{equation}

(2.7) 的左边是扰动 $f$ 所得的 Lefschetz 纤维丛的 Fukaya–Seidel 范畴的导出范畴，右边是在适当分次之下 $\check{f}$ 的矩阵因子化（matrix factorization）的范畴。矩阵因子化是 [Eis80] 在研究完全交上的 Cohen–Macaulay 模时引入的概念，在 Cohen–Macaulay 模的稳定范畴（stable category）的研究中被有效地使用\footnote{关于矩阵因子化的教科书，例如有 [Yos90]。}。矩阵因子化的同伦范畴等价于有界导出范畴（bounded derived category）的完全导出范畴（perfect derived category）所作的 Verdier 商 [Buc87]。Landau–Ginzburg 模型中 B 型膜由矩阵因子化给出这一想法归功于 Kontsevich\footnote{例如参见 [KL03]。}。对单纯奇点，已知猜想 2.7 的左边由 [Sei01a, Proposition 3.4]、右边由 [KST07] 分别与对应 Dynkin 箭图的表示的导出范畴等价，由此猜想 2.7 得证。更一般地，对于作为 A 型或 D 型多项式的 Sebastian–Thom 和所得的任意可逆多项式，猜想 2.7 也成立，这已由 [FU11, FU13] 证明。

% [Tutor Note: 原文此处作「Sebastian–Thom 和」，通行的名称为 Thom–Sebastiani 和（日文常作トム–セバスティアーニ和），原文疑为笔误或异写；译文照录，请审校确认。]

另外，在猜想 2.7 中，若把 $f$ 取为尖点奇点（cusp singularity）的定义多项式 $x^p+y^q+z^r+xyz$ 而非可逆多项式，则其镜像不是转置多项式，而是 $\PP^1$ 上 3 点处分别具有阶为 $p$、$q$ 及 $r$ 的固定子群的光滑有理轨形 $\PP^1_{a,b,c}$，这由 [Tak10] 猜想、[Kea15] 证明。这里 $(p,q,r)$ 是满足 $1/p+1/q+1/r<1$ 的正整数三元组。这与 [Ros10, TS15] 及 [Sei11] 密切相关。另外，把这里的辛几何侧与复几何侧互换所得之物在 [CHL] 中被研究。

关于奇点的镜像对称，也有不取定义多项式所定的 Lefschetz 纤维丛的 Fukaya–Seidel 范畴、而取变形奇点所得的 Liouville 区域上的卷绕 Fukaya 范畴（wrapped Fukaya category）的流派。对单纯奇点参见 [CU13, EL]，对尖点奇点参见 [Kea]，对 3 维通常二重点参见 [CPU]。

关于镜像对称最著名的教科书是 [CK99]。另外，2000 年与 2002 年 Clay 数学研究所主办的关于镜像对称的春季学校的记录 [HKK$^{+}$03, ABC$^{+}$09] 也可参考。关于镜像对称的综述有 [Hos99]，关于相干层的导出范畴的综述有 [Kaw06a]，关于 Lagrange 相交 Floer 理论的综述有 [Ono06, Fuk11]。关于同调镜像对称的综述有 [Fuk01, Fuk02a] 及 [Bal08]、[Por]。关于镜像对称与热带几何的关系参见 [Gro11]，关于 SYZ 猜想参见 [Gro13] 及其参考文献。另外，与 [KS06] 和 [GS11] 相关联，镜像对称与簇代数之间存在深刻关系这一点正逐渐明朗 [KS, GS15, GHKK, STWZ]。

\section{导出范畴}

同调镜像对称提供了把增广三角范畴看作空间、研究其几何学这一崭新的观点。也就是说，增广三角范畴是几何实体；至于它实现为某个流形的相干层的导出范畴、实现为另一个流形的 Fukaya 范畴的导出范畴、还是借助 Landau–Ginzburg 模型来实现，就像给流形安放坐标一样，并非本质。

Abel 范畴（Abelian category）的导出范畴（derived category），是把复形的同伦范畴用拟同构（quasi-isomorphism）局部化而得到的\footnote{范畴的局部化是环的局部化的自然推广。注意，只有一个对象的加法范畴是与环等价的概念。指这一点，范畴也被称为具有多个对象的环（rings with several objects）[Mit72]。沿着这一哲学，不仅局部化，模、Hochschild 上同调等关于环的许多概念都被推广到范畴。}。若 Abel 范畴具有足够多的内射对象（have enough injectives），则其导出范畴等价于由内射对象（injective object）构成的复形的同伦范畴。导出范畴常常用 [Ver96]\footnote{这是 Verdier 1967 年的学位论文，经若干编辑后于其去世后的 1996 年出版。}所引入的三角范畴（triangulated category）的语言来记述，但由于映射锥（mapping cone）缺乏函子性\footnote{若三角范畴具有函子性的映射锥，则它成为半单 Abel 范畴的充实子范畴 [Ver96, Proposition 1.2.13]。}、存在作为三角范畴等价但 Quillen 高阶 K 群不同构的 Abel 范畴的例子、以及仅凭三角范畴的结构不能从生成元复原导出范畴等各种理由，它被认为是把导出范畴定式化时并不令人满意的框架。取上同调会丢失信息，因此保持复形原样处理，是导出范畴的哲学。在三角范畴中，对象虽提升为复形，态射却取上同调（即只考虑复形之间态射的映射像，把同伦等价的态射视为同一），这一点是各种不便的原因。作为解决这一点的框架，有微分分次范畴、A∞ 范畴、稳定 ∞-范畴等；已知在特征 0 的域上它们给出等价的理论\footnote{微分分次范畴是 A∞ 范畴的特殊情形。另外，与给定 A∞ 范畴等价的微分分次范畴，可利用 A∞ 范畴中的米田引理来构造 [KS09, Section 6]。微分分次范畴与稳定 ∞-范畴理论的等价性参见 [Coh]。}。

在这些理论中，引入时所需预备知识最少的是微分分次范畴（differential graded category）。它是在通常的范畴公理中，把态射空间从单纯的线性空间提升为线性空间的复形，并对态射的合成施加 Leibniz 法则而成之物。

微分分次范畴的典型例子是如下定义的复形范畴。设 $\cA$ 为加法范畴。$\cA$ 的复形是指：$\cA$ 的对象列 $X := (X^i)_{i\in\Z}$ 与态射列 $d_X := (d_X^i)_{i\in\Z}$（其中 $d_X^i\in\Hom_{\cA}(X^i,X^{i+1})$）组成的对 $X^{\bullet} := (X,d_X)$，且对任意 $i\in\Z$ 满足 $d_X^{i+1}\circ d_X^i=0$。对复形 $X^{\bullet}$ 与 $Y^{\bullet}$，分次向量空间 $\hom(X,Y) := \bigoplus_{i\in\Z}\hom^i(X^{\bullet},Y^{\bullet})$ 由 $\hom^i(X^{\bullet},Y^{\bullet}) := \prod_{j\in\Z}\Hom_{\cA}(X^j,Y^{i+j})$ 定义。对 $\phi=(\phi^j\colon X^j\to Y^{j+i})_{j\in\Z}\in\hom^i(X^{\bullet},Y^{\bullet})$，定义映射 $d=(d^i\colon\hom^i(X^{\bullet},Y^{\bullet})\to\hom^{i+1}(X^{\bullet},Y^{\bullet}))_{i\in\Z}$ 为
\begin{equation}
    d^i\phi := \Bigl( (d^i\phi)^j := d_Y^{j+i}\circ\phi^j - (-1)^i\phi^{j+1}\circ d_X^j \colon X^j\to Y^{j+i+1} \Bigr)_{j\in\Z}\in\hom^{i+1}(X^{\bullet},Y^{\bullet}) \tag{3.1}
\end{equation}
则容易看出它成为微分（即满足 $d^{i+1}\circ d^i=0$）。态射的合成，对 $\phi\in\hom^i(X^{\bullet},Y^{\bullet})$ 与 $\psi\in\hom^j(Y^{\bullet},Z^{\bullet})$，由
\begin{equation}
    \psi\circ\phi := \Bigl( \psi^{k+i}\circ\phi^k \colon X^k\to Z^{i+k} \Bigr)_{k\in\Z}\in\hom^{i+j}(X^{\bullet},Z^{\bullet}) \tag{3.2}
\end{equation}
定义。态射的微分与合成满足 Leibniz 法则
\begin{equation}
    d(\psi\circ\phi)=(d\psi)\circ\phi+(-1)^j\psi\circ(d\phi) \tag{3.3}
\end{equation}
这一点很容易看出。

% [Tutor Note: 原文 (3.2) 中写作「$X^k\to Z^{i+k}$」，按 $\phi\in\hom^i(X,Y)$、$\psi\in\hom^j(Y,Z)$ 及本式右端 $\hom^{i+j}(X,Z)$ 的指标，似应为 $Z^{i+j+k}$；译文照原文保留，请审校确认。]

给定微分分次范畴 $\cD$ 时，可以不变更对象而取态射空间 0 次上同调来作加法范畴，称之为 $\cD$ 的上同调范畴（cohomology category），记作 $H^0(\cD)$。

任意态射都具有映射锥的微分分次范畴称为前三角（pretriangulated）。利用扭复形（twisted complex），任意微分分次范畴都可作为充实子范畴嵌入前三角微分分次范畴 [BK90]。前三角微分分次范畴的同伦范畴自然具有三角范畴的结构。对给定三角范畴 $\cT$，称同伦范畴与 $\cT$ 作为三角范畴等价的微分分次范畴（或 A∞ 范畴、或稳定 ∞-范畴）为 $\cT$ 的增广（enhancement）。把导出范畴作为增广三角范畴处理，上述 3 个困难便全部消解（例如参见 [Kel06, Toë11]）。

A∞ 范畴（A∞-category）是态射的合成在相差相容同伦的意义下结合的（associative up to coherent homotopy）这样的微分分次范畴的推广\footnote{A∞ 范畴的入门性文献，例如有 [Kel01, Kel02]。}。在相差同伦的意义下结合，与在相差相容同伦的意义下结合，其差别正是从微分分次范畴移到其同伦范畴时所丢失的信息。在 A∞ 范畴中，通过考虑态射的高次合成，可以恢复这一信息。A∞ 范畴除态射的微分与合成之外，一般还具有无限多个高次合成；态射的合成不满足结合律，但在上同调的层面上满足结合律，而且与结合律的偏离由态射的高次合成具体地描述出来。对微分分次范畴的前三角性\footnote{关于前三角 A∞ 范畴，例如参见 [BLM08] 及其参考文献。}、扭复形等概念自然推广到 A∞ 范畴，用它们定义 A∞ 范畴的导出范畴。

微分分次范畴构成的范畴 $\mathrm{dgcat}$ 又成为微分分次范畴，A∞ 范畴构成的范畴 $A_\infty\text{-}\mathrm{cat}$ 又成为 A∞ 范畴。高次合成全为零的 A∞ 范畴与微分分次范畴是同一概念，因此 $\mathrm{dgcat}$ 是 $A_\infty\text{-}\mathrm{cat}$ 的部分范畴；但两个微分分次范畴之间的 $A_\infty\text{-}\mathrm{cat}$ 中的态射比 $\mathrm{dgcat}$ 中的态射多得多，所以嵌入 $\mathrm{dgcat}\to A_\infty\text{-}\mathrm{cat}$ 不是充实的（full）。

$\mathrm{dgcat}$ 中使 $\mathrm{dgcat}$ 中的微分消失的态射称为微分分次函子（differential graded functor），$A_\infty\text{-}\mathrm{cat}$ 中使 $A_\infty\text{-}\mathrm{cat}$ 中的微分消失的态射称为 A∞ 函子（A∞-functor）。另外，诱导同伦范畴上等价的微分分次函子或 A∞ 函子称为拟等价（quasi-equivalence）。在 $H^0(\mathrm{dgcat})$ 中拟等价的逆函子未必存在，但在 $H^0(A_\infty\text{-}\mathrm{cat})$ 中任意拟等价都有逆函子存在。这一事实是与微分分次范畴相比 A∞ 范畴的优点之一。

微分为零的 A∞ 范畴称为极小模型（minimal model）\footnote{A∞ 范畴的极小模型与 Sullivan 的 de Rham 同伦理论中微分分次环的极小模型是不同的概念。}。对任意 A∞ 范畴，在 $A_\infty\text{-}\mathrm{cat}$ 中存在与之拟等价的微分分次范畴与极小模型。在微分分次范畴的范围内，即使任意对象之间的态射空间的上同调是有限维的，也未必存在拟等价的、任意对象之间态射空间为有限维的微分分次范畴。在 A∞ 范畴中，通过取极小模型，可以构成拟等价的、态射空间为有限维的 A∞ 范畴。这一事实也是与微分分次范畴相比 A∞ 范畴的优点之一。高次合成是为这些优点所必须付出的代价：$\hom(X,Y)$ 的有限维性，是用无限多个高次合成换来的。

拟等价的逆函子与极小模型，是用同调扰动理论（homological perturbation theory）构成的。关于同调扰动理论，例如参见 [Hue] 及其参考文献。

微分分次范畴与极小模型是 A∞ 范畴中的两极；当二者同时被满足时，即微分与高次合成两者都为零时，A∞ 范畴就只是分次范畴（graded category）。与分次范畴拟等价的 A∞ 范畴称为形式的（formal）。

奇点论中有称为有限决定性（finite determinacy）的概念，指全纯函数的孤立临界点的同构类由足够高阶的 jet 决定 [Mat68, Tou68]。与此类似，当 A∞ 范畴的拟等价类由直到足够高次合成为止的信息所决定时，称该 A∞ 范畴具有有限决定性（finitely determined）。有限决定性的关键是 Hochschild 上同调\footnote{这是与奇点论中的 Jacobi 环（也称 Milnor 环）相对应的。}的有限维性。

A∞ 结构的概念由 [Sta63] 引入，用于环路空间的特征化。据此，CW 复形 $X$ 能够环路空间化（deloop）\footnote{即存在某个 $Y$，使 $X$ 与基于点的环路空间 $\Omega Y$ 同伦等价。}的必要充分条件是 $X$ 为 A∞ 空间。A∞ 空间若用 [May72] 中引入的 operad 概念，则可定义为由 Stasheff 结合多面体（Stasheff associahedra）构成的拓扑 operad（topological operad）上的代数（algebra）。取这一拓扑 operad 的最高次链，就得到标准微分分次 A∞ operad（the standard differential graded A∞-operad）。标准微分分次 A∞ operad 上的代数就是 A∞ 代数。标准微分分次 A∞ operad 也可理解为结合 operad（associative operad）的自由分解（free resolution）[Mar96]。

关于 operad，结合代数的 Hochschild 复形成为小圆盘 operad（little disk operad）上的代数这一 Deligne 猜想（Deligne conjecture）很重要。关于 Deligne 猜想的证明参见 [Tam]，关于其推广参见 [Lur, Section 5.3] 及其参考文献。这与 [Kon03] 所证明并应用于 Poisson 流形形变量子化的 Kontsevich 形式性（Kontsevich formality）密切相关 [KS00, DTT09]。

有限决定性与 Kontsevich 形式性被有效地用于亏格 2 曲线的同调镜像对称的证明 [Sei11]。

A∞ 范畴是基于同伦论想法对微分分次范畴作出的灵活而精致的推广，而稳定 ∞-范畴（stable ∞-category）则是把同伦论进一步推到前台的导出范畴的定式化\footnote{关于 ∞-范畴参见 [Lur09]，关于稳定 ∞-范畴参见 [Lur] 及其参考文献。}。∞-范畴（∞-category）也称为 $(1,\infty)$-范畴（$(1,\infty)$-category），指高次态射全部可逆的高次范畴（higher category）。∞-范畴有拓扑范畴（topological category）、单纯增广范畴（simplicially enriched category）、Segal 范畴（Segal category）、拟范畴（quasicategory）等各种模型。取 ∞-范畴的同伦范畴，可以得到去掉高次信息的通常意义的范畴。范畴的局部化是 ∞-范畴出现的典型情形，把复形的同伦范畴用拟同构局部化所得的导出范畴即属此类。

给定范畴 $\cC$ 与其上称为弱等价（weak equivalence）的态射族 $\cW$ 时\footnote{范畴 $\cC$ 上的弱等价，是指 $\cC$ 的部分范畴 $\cW$，它包含所有同构（特别地，$\cC$ 的所有对象都是 $\cW$ 的对象），并满足 2-out-of-3 性（two-out-of-three property，即对任意图表 $X\xrightarrow{f}Y\xrightarrow{g}Z$，若 $f,g,g\circ f$ 中有两个属于 $\cW$，则剩下的一个也属于 $\cW$）。}，$\cC$ 关于 $\cW$ 的局部化（localization）$\cC[\cW^{-1}]$，定义为这样的范畴：使 $\cW$ 的元素变为同构的、从 $\cC$ 出发的函子，都唯一地经由它。Dwyer–Kan [DK80] 的单纯局部化（simplicial localization）是由组 $(\cC,\cW)$ 给出单纯增广范畴 $L(\cC,\cW)$ 的方法；取这一单纯增广范畴的同伦范畴，便得到通常的局部化：$\cC[\cW^{-1}]\simeq\operatorname{Ho}L(\cC,\cW)$。局部化后的范畴 $\cC[\cW^{-1}]$ 与 $\cC$ 有相同的对象，但具体描述 $\cC[\cW^{-1}]$ 中的态射一般很困难。若弱等价 $\cW$ 具有分数计算（calculus of fraction），情况会稍有改善，但本质的进步是由模型范畴（model category）带来的。这是作为同伦论中构造的抽象化、由 Quillen 引入的概念，是同调代数在某种意义上的非线性化。模型范畴除弱等价 $\cW$ 之外，还具有纤维化（fibration）、余纤维化（cofibration）以及称为函子性因子化（functorial factorization）的额外结构，并且这些结构被施加了强公理，使得局部化范畴中的态射能够被具体描述\footnote{对给定的范畴与弱等价的组，在其上添加纤维化与余纤维化以构成模型，一般并非显然的问题。另外，即便可能，纤维化与余纤维化的添加方式也不唯一。}。

这些公理包含完备性（completeness）与余完备性（cocompleteness），即任意极限（limit）及余极限（colimit）的存在；特别地，作为空图表（empty diagram）的余极限与极限，定义了始对象（initial object）与终对象（terminal object）。从始对象到终对象有唯一的态射，当它是同构时范畴称为带点的（pointed）。对不带点的范畴 $\cC$，通过考虑终对象 $\ast$ 的下范畴\footnote{对范畴 $\cC$ 及 $\cC$ 的对象 $S$，$S$ 的下范畴（under category 或 coslice category）是指以 $\cC$ 的态射 $S\to X$ 为对象、以 $S$ 之下的交换图表为态射的范畴 $\cC_{S/}$。} $\cC_{\ast}:=\cC_{\ast/}$，可以标准地构成带点范畴\footnote{位相空间的范畴具有终对象，但不具有始对象，所以不是带点范畴。另一方面，带点位相空间构成的范畴是带点范畴。}。

% [Tutor Note: 原文此处写作「終対象 ∗ の下圏 $\mathcal{C}_* := \mathcal{C}_{*/}$」，等号左右的下标不一致（下范畴的通常记号应为 $\mathcal{C}_{*/}$），疑为笔误；译文照原文保留，请审校确认。]

带点 ∞-范畴 $\cC$ 的三角（triangle）是指
\begin{equation}
    \begin{array}{ccc}
        X & \xrightarrow{\hspace*{7mm}f\hspace*{7mm}} & Y\\[9pt]
        \downarrow & & \downarrow\scriptstyle{g}\\[9pt]
        \ast & \xrightarrow{\hspace*{19mm}} & Z
    \end{array} \tag{3.4}
\end{equation}
形式的图表。(3.4) 是拉回图表时，称 $X\xrightarrow{f}Y\xrightarrow{g}Z$ 为完全三角（exact triangle），并称 $f$ 为 $g$ 的核（kernel）。与此对偶地，(3.4) 是推出图表时，称 $X\xrightarrow{f}Y\xrightarrow{g}Z$ 为余完全三角（coexact triangle），并称 $g$ 为 $f$ 的余核（cokernel）。任意态射都具有核与余核，且完全三角与余完全三角一致的带点 ∞-范畴就是稳定 ∞-范畴（stable ∞-category）。这与由 $\Sigma X := \ast\coprod_{X}\ast$ 确定的悬垂函子 $\Sigma$（或由 $\Omega X := \ast\times_{X}\ast$ 确定的环路函子 $\Omega$）成为 ∞-范畴的自等价是等价的，「稳定」一词即由此而来。对稳定 ∞-范畴的对象 $X$，
\begin{equation}
    X[n] :=
    \begin{cases}
        \Sigma^n(X) & n\ge0,\\
        \Omega^{-n}(X) & n<0
    \end{cases} \tag{3.5}
\end{equation}
写。

稳定 ∞-范畴的同伦范畴成为 Verdier 意义的三角范畴\footnote{例如参见 [Lur]。另外参照 [Hov99, Chapter 7]。}。与三角范畴比较时，稳定 ∞-范畴的特征之一在于：稳定 ∞-范畴是在 ∞-范畴上添加条件之物，而不是添加结构之物。

∞-范畴之间的函子与有限极限交换时称为左正合的（left exact）。同样地，∞-范畴之间的函子与有限余极限交换时称为右正合的（right exact）。左正合且右正合的函子称为正合的（exact）。对稳定 ∞-范畴

% 本 chunk 依原文在 p009 页末截断（原文「安定 ∞ 圏に対」）：c02 自 p010「しては…」接续。
% CHUNK c01 DONE pages p006–p009
